\documentclass[11pt]{article}
\usepackage[margin=25mm]{geometry}
\usepackage{fontspec}
\setmainfont{Latin Modern Roman}
\usepackage{amsmath,amssymb,mathtools}
\usepackage{unicode-math}
\setmathfont{Latin Modern Math}
\usepackage{xurl,hyperref}
\hypersetup{colorlinks=true,urlcolor=blue}
\setcounter{secnumdepth}{0}
\setlength{\parindent}{0pt}
\setlength{\parskip}{5pt}
\setlength{\emergencystretch}{3em}
\newcommand{\sref}[2]{\hyperlink{source:#1:#2}{[S#2]}}
\newcommand{\cataloguescope}[1]{\paragraph{Recorded target.}#1}
\newcommand{\R}{\mathbb R}
\newcommand{\C}{\mathbb C}
\newcommand{\N}{\mathbb N}
\newcommand{\Z}{\mathbb Z}
\newcommand{\Q}{\mathbb Q}
\newcommand{\F}{\mathbb F}
\newcommand{\eps}{\varepsilon}
\begin{document}
\section{Self-Consistent Motion of Massive Charged Particles in Relativity}\label{statement}

\subsection{Definitions and mathematical statement}

\paragraph{Shared notation.}$\mathbb N=\{1,2,\ldots\}$; the standard number fields have their usual meanings. All additional conventions are stated locally.


Massive particles follow timelike worldlines. In the singular-particle description these curves are singularities of a Lorentzian metric and an electromagnetic field. Away from them the fields satisfy the Einstein--Maxwell equations; their distributional extension and conservation laws must be defined rather than inferred from smooth identities.

\paragraph{Catalogue wording retained for traceability.}
How can one rigorously model the motion of massive and charged particles in general relativity, including their coupled electromagnetic and gravitational fields and the timelike singularities they generate?

\paragraph{Scope and source qualification.}The report asks for a modelling framework, not a uniquely specified existence theorem. A distributional class, regularization and notion of particle motion must be stated in any proposed solution.
\subsection{Short English statement}
Develop a mathematically consistent coupled description of massive charged particles and their electromagnetic and gravitational fields.
\subsection{Research attempt}
\paragraph{Model, date and investigated scope.}
OpenAI GPT-6 Astra (configured model \texttt{gpt-6-astra}), ultra reasoning; generated 11 October 2026. Computed the nonintegrable Coulomb self-energy and located the undefined stress/self-force step in a naive point-particle framework.

\paragraph{Sources and literature check.}
Read the supplied definition, catalogue wording, source qualifications and reference entries for this question. The historical source/context and inherited status are retained; no fresh literature or novelty certification is claimed. No external theorem is used without proof in this attempt; explicit model assumptions are distinguished from facts about the source model.

\paragraph{The point-charge energy obstruction.} In flat three-space, choose electromagnetic units with Coulomb field \(E(x)=q x/(4\pi|x|^3)\), \(q\ne0\), and energy density \(|E|^2/2\). For a punctured ball,
\[
 \int_{\varepsilon<|x|<R}\frac{|E|^2}{2}\,dx
 =\frac{q^2}{8\pi}\int_\varepsilon^R\frac{dr}{r^2}
 =\frac{q^2}{8\pi}\left(\frac1\varepsilon-\frac1R\right).
\]
Thus the ordinary quadratic electromagnetic stress is not locally integrable at the worldline when integrated over space and any nonzero finite time interval. The field itself is locally integrable in space, since its magnitude scales as \(r^{-2}\), but squaring it destroys that property.
\paragraph{Failure of a distributional shortcut.} Writing Maxwell's equation with a delta current does not automatically define the stress tensor needed in Einstein's equation, nor the Lorentz force obtained by evaluating the self-field on the same curve. A cutoff produces the explicit divergent mass-scale term above, so removing it requires a specified renormalization and a proof that remaining quantities and conservation laws are well defined.
\paragraph{Source gap.} Curved geometry or nonlinear electromagnetic models may change the local behavior; this flat calculation is not an impossibility theorem for the source programme. It does show why a naive distributional point-particle substitution is incomplete. The first gap is a self-consistent singularity/regularization class in which both gravitational and electromagnetic stress and particle motion exist with controlled limiting conservation. No such coupled existence theorem or trajectory law is constructed.

\paragraph{Verification and adversarial audit.}
Checked radial constants and powers, local integrability of the field versus its square and the limitation of a flat-space diagnostic. This is a GPT-6 Astra Ultra self-audit, not independent human verification. Prompt A drafting followed by prompt B scope/dependency/boundary self-audit. Checked radial constants and powers, local integrability of the field versus its square and the limitation of a flat-space diagnostic. The remaining gap is: Construct a singularity or renormalization framework defining coupled stress and self-force and proving conservation and motion in the limit. Unresolved status and zero original-target completion are retained.

Coordinator checked the Coulomb energy coefficient and local integrability distinction between the field and its quadratic stress. This additional check was model-assisted, not independent human verification.

\paragraph{Tools, process and contributions.}
Python 3 and local file-reading/writing tools for manuscript serialization and structural checks; mathematical calculations were symbolic reasoning, with no numerical experiment or theorem prover. Contribution prompt A was followed by prompt B. GPT-6 Astra with ultra reasoning drafted the research and performed its stated self-audit; any additional coordinator contribution is identified in the verification and audit. A GPT-6-based Codex coordinator assembled the manuscript and checked repository integration; its exact variant and reasoning configuration were not exposed. No human mathematical contribution, independent human proof check or formal proof certification is supplied. Complete derivations and the consulted source references are included here; private reasoning transcripts are not published.

\paragraph{Final status and first remaining gap.}
\textbf{UNRESOLVED}. Construct a singularity or renormalization framework defining coupled stress and self-force and proving conservation and motion in the limit.

\section{Completion Estimate}
\noindent\textbf{COMPLETION: 0\%}
The proved calculation is an elementary auxiliary result or a test of a restricted strategy; no new portion of the full source target is established.

\subsection{Sources}
\begin{itemize}
\item[\hypertarget{source:statement:1}{S1}] ULAM.AI and the UnsolvedMath Contributors, UnsolvedMath, catalogue entry. Catalogue curation under CC BY 4.0; definitions and English exposition added locally.
\url{https://huggingface.co/datasets/ulamai/UnsolvedMath}.
\item[\hypertarget{source:statement:2}{S2}] Mathematical Aspects of General Relativity (2021), Oberwolfach Reports, DOI 10.4171/owr/2021/40, pp. 2157--2267. Original report and mathematical context.
\url{https://doi.org/10.4171/owr/2021/40}.
\end{itemize}
\label{end:statement}
\end{document}
